% Teil II: der durchgerechnete Weg.

\subsection*{Gegeben und gesucht}
Gegeben: $M=\qty{2}{kg}$, $m=\qty{10}{g}=\qty{0.010}{kg}$,
$L=\qty{1}{m}$, $\varphi=37^\circ$, $g=\qty{9.81}{m/s^2}$.
Gesucht: (a) $v_\mathrm{Kugel}$; (b) Deformationsenergie; (c)
Kugelgeschwindigkeit für Überschlag (Stange bzw. Seil); (d) Geschwindigkeit
des harten Klotzes $V$ nach elastischem Rückprall der Kugel mit
$c=\qty{10}{m/s}$.

\subsection*{Ansatz}
(a) Aus $\varphi$ die Hubhöhe $h$; Energie während des Pendelumlaufs, dann
Impuls am Stoss. (b) Differenz der kinetischen Energien.
(c) Überschlag: Energiebetrachtung bis zum höchsten Punkt. (d)
Elastischer Stoss: Impuls- und Energieerhaltung.

\subsection*{Rechnung}
\textbf{a) Kugelgeschwindigkeit.}
\begin{equation*}
  h = L(1-\cos\varphi)=\qty{1}{m}(1-\cos 37^\circ)=\qty{0.201}{m}.
\end{equation*}
Energie nach dem Stoss (gemeinsame Geschwindigkeit $u$):
\begin{equation}
  \frac12 (m+M)u^2 = (m+M)gh
  \;\Rightarrow\;
  u=\sqrt{2gh}=\qty{1.987}{m/s}.
\end{equation}
Impuls beim Stoss, $m v=(m+M)u$:
\begin{equation}
  v = \frac{m+M}{m}\,u
  = \frac{2.010}{0.010}\,\qty{1.987}{m/s}
  \approx \qty{400}{m/s}.
\end{equation}

\textbf{b) Deformationsenergie:}
\begin{equation}
  \Delta E = \frac12 m v^2 - \frac12 (m+M)u^2
  = \frac12\cdot 0.01\cdot 400^2\ \si{J}
    - \frac12\cdot 2.01\cdot 1.987^2\ \si{J}
  \approx \qty{794}{J}.
\end{equation}

\textbf{c) Überschlag.} Im höchsten Punkt ($h=2L$) muss für die Stange nur
$u_{\mathrm{top}}\ge 0$ sein; Energie:
\begin{equation}
  \frac12 (m+M)u^2 = (m+M)\,g\cdot 2L\;\Rightarrow\;
  u=\sqrt{4gL}=\qty{6.26}{m/s}
  \;\Rightarrow\;
  v = \frac{m+M}{m}u=\qty{1259}{m/s}.
\end{equation}
Beim \emph{Seil} muss im obersten Punkt $u_{\mathrm{top}}^2/L \ge g$ gelten
(für $u_{\mathrm{top}}=\sqrt{gL}$), also
\begin{equation}
  \frac12 (m+M)u^2 = (m+M)g\cdot 2L + \frac12(m+M)gL
  \;\Rightarrow\; u=\sqrt{5gL}=\qty{7.00}{m/s}
  \;\Rightarrow\; v=\qty{1408}{m/s}.
\end{equation}
Mit dem Seil ist also mehr Kugelgeschwindigkeit nötig: Bei
$\varphi=180^\circ$ reicht die zur Verfügung stehende Energie zum Erreichen
der Höhe $2L$ nicht mehr mit der Zusatzbedingung der Zentripetalbeschleunigung
aus. Antwort: Mit der Stange $\qty{1259}{m/s}$, mit dem Seil $\qty{1408}{m/s}$
nötig - die erste reicht für das Seil \emph{nicht}.

\textbf{d) Klotzgeschwindigkeit nach elastischem Rückstoss.} Impuls- und
Energieerhaltung mit Kugelgeschwindigkeiten vor ($v$) und nach ($-c$), Klotz
$V$:
\begin{equation*}
  m v = m(-c) + M V\;,\qquad
  \frac12 m v^2 = \frac12 m c^2 + \frac12 M V^2.
\end{equation*}
Aus der zweiten: $v^2 = c^2 + \frac{M}{m}V^2$; zusammen mit der ersten
$v = -c+\frac{M}{m}V$ führt das (Lösung nach $V$) auf
\begin{equation}
  V = \frac{2mc}{M-m}
  = \frac{2\cdot 0.010\cdot 10}{2-0.010}\ \si{m/s}
  = \qty{0.101}{m/s}.
\end{equation}

\begin{ergebnis}
(a) $v\approx\qty{400}{m/s}$; (b) $\Delta E\approx\qty{794}{J}$;
(c) Stange $v=\qty{1.26e3}{m/s}$, Seil $v=\qty{1.41e3}{m/s}$ - mit dem
Seil erreicht man die Höhe $2L$ mit $v=\qty{1259}{m/s}$ nicht;
(d) $V=\dfrac{2mc}{M-m}=\qty{0.101}{m/s}$.
\end{ergebnis}

\begin{kontrolle}
\textbf{Einheiten:} In (a)
$v=\tfrac{m+M}{m}\sqrt{2gh}$: $kg/kg\cdot\sqrt{m^2/s^2}=m/s$. \checkmark
Grenzfall (c): ein Überschlag benötigt bei einem Seil mehr Energie als bei
einer Stange, weil am obersten Punkt zusätzliche Zentripetalbeschleunigung
bereitstehen muss - plausibel.

\textbf{Elastischer Limit:} Für $M\to m$ (in d) teilen sich Energie und
Impuls; $\frac{2m c}{M-m}$ wird für $M\to m$ gross (Kugel wird fast nicht
gebremst). Einheiten: $kg\cdot m/s/kg = m/s$. \checkmark
\end{kontrolle}

\begin{merksatz}
Beim vollständig inelastischen Stoss ist nur der Impuls erhalten, bei
elastischem zusätzlich die kinetische Energie. Beim ballistischen Pendel
koppeln Stoss (Impulserhaltung) und Pendel (Energieerhaltung) über die
gemeinsame Geschwindigkeit $u$ direkt nach dem Stoss.
\end{merksatz}